\section{PTA 刷题总结}
%\textcolor{blue}{\textbf{271A}}
%\begin{lstlisting}[language=C++]\end{lstlisting}
在本节中，我将总结我在 PTA 上做题的经验。

\subsection{遇到的问题以及学到的方法}
\begin{itemize}
\subsection{对数组每一列按升序排序}
\item 题目链接：\href{https://pintia.cn/problem-sets/1997490617200599040/exam/problems/type/7?problemSetProblemId=1997490617559937029}{坚固矩阵(点击此处)}；						

代码示例：\\
\begin{lstlisting}[language=C++]
void solve() {
	ll n, m;
	cin >> n >> m;
	vector<vector<ll> > v(n, vector<ll>(m));
	for (ll i = 0; i < n; i++) {
		for (ll j = 0; j < m; j++) {
			cin >> v[i][j];
		}
	}
	for (ll j = 0; j < m; j++) {
		//开一个新数组来存每一列的数字，然后排序；
		vector<ll> a;
		for (ll i = 0; i < n; i++) {
			a.push_back(v[i][j]);
		}
		sort(a.begin(), a.end());
		//将排序好的数组在放进原数组中；
		for (ll i = 0; i < n; i++) {
			v[i][j] = a[i];
		}
	}
	bool flag;
	for (ll i = 0; i < n; i++) {
		flag = true;
		for (ll j = 0; j < m; j++) {
			cout << (flag == true ? "" : " ") << v[i][j];
			flag = false;
		}
		cout << endl;
	}
}
\end{lstlisting}
\subsection{轮流问题}
\item 题目链接：\href{https://pintia.cn/problem-sets/1997490617200599040/exam/problems/type/7?problemSetProblemId=1997490617559937026}{小杭玩游戏(点击此处)}；总结：可以用bool来表示轮到谁了。

代码示例：\\
\begin{lstlisting}[language=C++]
void solve() {
	ll a,b,n,cnt=0;
	cin>>a>>b>>n;
	bool flag=true;
	while(n>0){
		ll m;
		cnt++;
		if(flag){
			m=gcd(n,a);
		}else{
			m=gcd(n,b);
		}
		n-=m;
		flag=!flag;
	}
	if(flag){
		cout<<"Xiao Hong Win"<<endl;
		cout<<cnt<<endl;
	}else{
		cout<<"Xiao Ming Win"<<endl;
		cout<<cnt<<endl;
	}
}
\end{lstlisting}
\subsection{结构体的使用}
\item 题目链接：\href{https://pintia.cn/problem-sets/994805046380707840/exam/problems/type/7?problemSetProblemId=994805071789801472&page=1}{L2-003月饼(点击此处)}；\\\textcolor{red}{总结：当要记录多个数据(两个以上)时可以用结构体来操作。}处理两个数据可以用哈希表或者其他容器操作。
\\lambda函数也要了解一下。

代码示例：\\
\begin{lstlisting}[language=C++]
class total {
	public:
	double w;
	double v;
	double d;
};

bool cmp(total t1, total t2) {
	return t1.d > t2.d;
}

void solve() {
	ll n;
	double m;
	cin >> n >> m;
	vector<total> a(n);
	for (ll i = 0; i < n; i++) {
		cin >> a[i].w;
	}
	for (ll i = 0; i < n; i++) {
		cin >> a[i].v;
		a[i].d = a[i].v / a[i].w;
	}
	sort(a.begin(), a.end(), cmp);
	//也可以写成lambda函数的形式
	//  sort(a.begin(),a.end(),[](total &t1,total &t2){
		//          return t1.d>t2.d;
		//  });
	double ans = 0;
	for (ll i = 0; i < n && m > 0; i++) {
		if (a[i].w <= m) {
			ans += a[i].v;
			m -= a[i].w;
		} else {
			ans += a[i].d * m;
			m = 0;
		}
	}
	cout << fixed << setprecision(2) << ans << endl;
}
\end{lstlisting}
\subsection{固定位数问题}
\item 题目链接：\href{https://pintia.cn/problem-sets/994805046380707840/exam/problems/type/7?problemSetProblemId=994805099426070528}{出生年(点击此处)}  总结：set去重可以检查有几个不同，for循环固定4次，就算temp=0也要继续循环，保证了始终有四位。

代码示例：\\
\begin{lstlisting}[language=C++]
int main() {
	int y, n;
	cin >> y >> n;
	for (int i = y; ; i++) {
		set<int> st;
		int temp = i;
		for (int j = 0; j < 4; j++) {
			st.insert(temp % 10);
			temp /= 10;
		}
		if (st.size() == n) {
			printf("%d %04d", i - y, i);
			return 0;
		}
	}
}
\end{lstlisting}
\subsubsection{字符串查找连续满足条件串的方法}
\item 题目链接：\href{https://pintia.cn/problem-sets/994805046380707840/exam/problems/type/7?problemSetProblemId=1111914599408664577}{L1-058 6翻了(点击此处)}；

代码示例：\\
\begin{lstlisting}[language=C++]
auto solve() -> void {
	string s;
	getline(cin, s);
	for (ll i = 0; i < s.size();) {
		if (s[i] == '6') {
			ll cnt = 0;
			ll j = i;
			//改成while循环代替之前的cnt=0和！=0的分类讨论
			//尽量用这种思想替换掉以取的思想
			while (j < s.size() && s[j] == '6') {
				cnt++;
				j++;
			}
			if (cnt > 9) {
				cout << "27";
			} else if (cnt > 3) {
				cout << "9";
			} else {
				for (ll k = 0; k < cnt; k++) {
					cout << "6";
				}
			}
			i = j;
		} else {
			cout << s[i];
			i++;
		}
	}
}
\end{lstlisting}
\subsubsection{活用字符串字串提取函数(substr)}
\item 题目链接：\href{https://pintia.cn/problem-sets/994805046380707840/exam/problems/type/7?problemSetProblemId=1111914599412858880}{L1-059 敲笨钟(点击此处)}；

代码示例：\\
\begin{lstlisting}[language=C++]
auto solve() -> void {
	string s;
	getline(cin, s);
	//要熟练的使用find和substr函数，对字符串的提取和替换都有很大的作用
	size_t pos1 = s.find(',');
	size_t pos2 = s.find('.');
	bool flag = false;
	if (pos1 >= 3 && pos2 >= 3) {
		if (s.substr(pos1 - 3, 3) == "ong" && s.substr(pos2 - 3, 3) == "ong") {
			flag = true;
		}
	}
	if (!flag) {
		cout << "Skipped" << endl;
	} else {
		ll i = s.size() - 1;
		ll cnt = 0;
		while (i > 0) {
			if (s[i] == ' ') {
				cnt++;
			}
			if (cnt == 3) {
				break;
			}
			i--;
		}
		//这种输出之前也没有见过，要学习一下，不要再一个一个遍历s再存t了
		cout << s.substr(0, i + 1) << "qiao ben zhong." << endl;
	}
}
\end{lstlisting}
\subsubsection{最小化思想2026天梯选拔赛1-6}
\item 题目链接：\href{https://pintia.cn/problem-sets/2030339464381566976/exam/problems/type/7?problemSetProblemId=2030339464444481541}{1-6源初之石(点击此处)}；

代码示例：\\
\begin{lstlisting}[language=C++]
auto solve() -> void {
	ll n;
	cin >> n;
	ll a = -1, b = -1, c = -1;
	//之所以要看i*i*i是为了找到能使三个数相乘=n的最小值(这里默认把最小值给了a)；
	for (ll i = 2; i * i * i <= n; i++) {
		if (n % i == 0) {
			a = i;
			break;
		}
	}
	if (a == -1) {
		cout << "NO" << endl;
		return;
	}
	//因为前面的a已经使最小的了所以这里找c的时候可以从a+1开始找，可以提高效率。
	for (ll i = a + 1; i * i <= n / a; i++) {
		if ((n / a) % i == 0) {
			c = i;
			b = n / a / c;
			if (b >= 2 && b != a && b != c) {
				cout << "YES" << endl;
				cout << a << " " << c << " " << b << endl;
				return;
			}
		}
	}
	cout << "NO" << endl;
}
\end{lstlisting}
\end{itemize}
\newpage


























